% Determinant-polynomial rigidities for the mixed-boundary discrete Laplacian L_k.

\begin{proposition}[行列式多项式的系数闭式：\texorpdfstring{$\binom{k+j}{2j}$}{binom(k+j,2j)} 公式与正系数刚性]\label{prop:pom-Lk-det-coeff-binomial}
对 \(k\ge 0\) 定义
\[
D_k(t):=\det(L_k+tI),
\qquad
D_0(t):=1.
\]
则 \(D_k(t)\) 为次数 \(k\) 的整系数首一多项式，且具有完全显式的系数闭式
\[
\boxed{\;
D_k(t)=\sum_{j=0}^{k}\binom{k+j}{2j}\,t^{j}\; }.
\]
特别地，所有系数严格为正，并满足端点刚性
\[
[t^0]D_k=1,\qquad [t^{k-1}]D_k=2k-1,\qquad [t^{k}]D_k=1.
\]
\leanverified{detPoly\_eval\_zero}
\leanverified{detPoly\_monic}
\leanverified{detPoly\_natDegree}
\leanverified{detPoly\_coeff\_sub\_leading}
\leanverified{detPoly\_deriv\_eval\_zero\_double}
\leanverified{detPoly\_coeff\_binomial}
\leanverified{detPoly\_coeff\_pos}
\leanverified{paper\_pom\_lk\_det\_coeff\_binomial}
\end{proposition}

\begin{proof}
由推论 \ref{cor:pom-Lk-shifted-det-free-energy} 的双曲形式（取 \(t>0\)）
\[
D_k(t)=\frac{\cosh((2k+1)\eta(t))}{\cosh(\eta(t))},
\qquad
\eta(t)=\mathrm{arsinh}\sqrt{\frac{t}{4}},
\]
并用恒等式
\(
\cosh((n+2)\eta)=2\cosh(2\eta)\cosh(n\eta)-\cosh((n-2)\eta)
\)
（\(n\ge 1\)），得到递推
\[
D_{k+1}(t)=(t+2)D_k(t)-D_{k-1}(t),
\qquad
D_0(t)=1,\quad D_1(t)=t+1,
\]
其中使用了 \(2\cosh(2\eta(t))=t+2\)。
由于两边均为关于 \(t\) 的多项式，上述递推对一切 \(t\in\CC\) 成立。
\par
令
\(
F_k(t):=\sum_{j=0}^{k}\binom{k+j}{2j}\,t^j
\)。
直接计算得 \(F_0=1,F_1=t+1\)。
并且对任意 \(j\ge 0\) 有系数恒等式
\[
\binom{k+1+j}{2j}
\;=\;
2\binom{k+j}{2j}+\binom{k+j-1}{2j-2}-\binom{k+j-1}{2j},
\]
其证明只需两次应用 Pascal 公式：
\[
\binom{k+1+j}{2j}
=\binom{k+j}{2j}+\binom{k+j}{2j-1}
=\binom{k+j-1}{2j}+2\binom{k+j-1}{2j-1}+\binom{k+j-1}{2j-2}.
\]
因此 \(F_{k+1}=(t+2)F_k-F_{k-1}\)。
由同一递推与初值的唯一性，得 \(D_k\equiv F_k\)。
端点系数与严格正性由闭式直接给出。证毕。
\end{proof}

\begin{corollary}[Lucas 型闭式与二次域单位表示：\texorpdfstring{$t(4+t)$}{t(4+t)} 判别式的结构通道]\label{cor:pom-Lk-det-lucas-unit}
\leanverified{paper\_pom\_lk\_det\_lucas\_unit}
记判别式 \(\Delta(t):=t(4+t)\)，并在 \(\CC\setminus[-4,0]\) 上固定 \(\sqrt{\Delta(t)}\) 的解析分支（与 \(t>0\) 的正实分支相容）。
定义
\[
r(t):=\frac{t+2+\sqrt{\Delta(t)}}{2},
\qquad
r(t)^{-1}=\frac{t+2-\sqrt{\Delta(t)}}{2}.
\]
则 \(r(t)r(t)^{-1}=1\)、\(r(t)+r(t)^{-1}=t+2\)，并且对一切 \(k\ge 0\) 有闭式
\[
\boxed{\;
D_k(t)=\frac{r(t)^{\,k+1}+r(t)^{-k}}{r(t)+1}\; }.
\]
从而 \((D_k(t))_{k\ge 0}\) 满足二阶 Lucas 递推
\[
D_{k+1}(t)=(t+2)D_k(t)-D_{k-1}(t),
\qquad
D_0(t)=1,\quad D_1(t)=t+1.
\]
\leanverified{detPoly\_zero}
\leanverified{detPoly\_one}
\leanverified{detPoly\_succ\_succ}
\leanverified{detPoly\_two}
\leanverified{detPoly\_three}
\leanverified{detPoly\_eval\_two\_zero}
\leanverified{detPoly\_eval\_two\_one}
\leanverified{detPoly\_eval\_two\_recurrence}
\leanverified{detPoly\_eval\_neg\_one\_rec}
\leanverified{detPoly\_eval\_neg\_one\_periodic}
\leanverified{detPoly\_eval\_neg\_one\_zero}
\leanverified{detPoly\_eval\_neg\_one\_one}
\leanverified{detPoly\_eval\_neg\_one\_two}
\leanverified{detPoly\_eval\_neg\_one\_three}
\leanverified{detPoly\_eval\_neg\_one\_four}
\leanverified{detPoly\_eval\_neg\_one\_five}
\end{corollary}

\begin{proof}
由推论 \ref{cor:pom-Lk-area-law-qform}，对 \(t>0\) 有
\(
D_k(t)=q(t)^{-k}\frac{1+q(t)^{2k+1}}{1+q(t)}
\)
且 \(q(t)=e^{-2\eta(t)}\in(0,1)\)。
令 \(r(t):=q(t)^{-1}=e^{2\eta(t)}\)，则
\[
D_k(t)=\frac{r(t)^{\,k+1}+r(t)^{-k}}{r(t)+1}.
\]
另一方面，由 \(r+r^{-1}=e^{2\eta}+e^{-2\eta}=2\cosh(2\eta)=t+2\) 得
\(
r(t)=\frac{t+2+\sqrt{t(4+t)}}{2}
\)
（取与 \(t>0\) 相容的根号分支），故得到所示表述。
Lucas 递推由 \(r+r^{-1}=t+2\) 的一阶消元直接得到。证毕。
\end{proof}

\begin{proposition}[Cassini--Pell 守恒律：行列式 Wronskian 恒等式]\label{prop:pom-Lk-det-cassini-pell}
对一切 \(k\ge 1\) 与一切复参数 \(t\in\CC\) 成立严格恒等式
\[
\boxed{\;
D_{k+1}(t)D_{k-1}(t)-D_k(t)^2=t.\;}
\]
该恒等式等价于把 \((D_{k-1}(t),D_k(t))\) 嵌入二次曲线
\(
X^2-(t+2)XY+Y^2=t
\)
的整点轨道；在 \(t=1\) 时退化为奇指标 Fibonacci 行列式的 Cassini 恒等式。
\leanverified{detPoly\_cassini\_pell}
\leanverified{paper\_pom\_lk\_det\_cassini\_pell}
\leanverified{detPoly\_eval\_one}
\leanverified{fib\_odd\_cassini}
\leanverified{detPoly\_eval\_cassini\_gap}
\end{proposition}

\begin{proposition}[统一 Fibonacci Cassini 恒等式]\label{prop:pom-fib-cassini-unified}
对所有 \(n \in \mathbb{N}\)，在 \(\mathbb{Z}\) 中成立
\[
F_n \cdot F_{n+2} = F_{n+1}^2 + (-1)^{n+1}.
\]
即：当 \(n\) 为奇数时 \(F_n F_{n+2} = F_{n+1}^2 + 1\)；当 \(n\) 为偶数时 \(F_n F_{n+2} = F_{n+1}^2 - 1\)。
\leanverified{fib\_product\_cassini}
\end{proposition}

\begin{proof}
由推论 \ref{cor:pom-Lk-det-lucas-unit} 的递推，定义
\(
W_k:=D_{k+1}D_{k-1}-D_k^2
\)。
直接代入 \(D_{k+1}=(t+2)D_k-D_{k-1}\) 得
\[
W_k=\bigl((t+2)D_k-D_{k-1}\bigr)D_{k-1}-D_k^2
=D_k\bigl((t+2)D_{k-1}-D_{k-2}\bigr)-D_{k-1}^2
=W_{k-1},
\]
故 \(W_k\) 与 \(k\) 无关。
取 \(k=1\) 得
\(
W_1=D_2D_0-D_1^2=((t+2)(t+1)-1)- (t+1)^2=t
\)，从而对所有 \(k\ge 1\) 有 \(W_k=t\)。证毕。
\end{proof}

\begin{corollary}[严格对数凸性与比值单调：Cassini--Pell 不变量的曲率锁定]\label{cor:pom-Lk-det-logconvex-ratio}
当 \(t>0\) 时，由 \(D_k(t)>0\)（\(L_k+tI\succ 0\)）与命题 \ref{prop:pom-Lk-det-cassini-pell} 得对一切 \(k\ge 1\) 成立严格不等式
\[
D_k(t)^2< D_{k-1}(t)\,D_{k+1}(t),
\qquad
\frac{D_{k-1}(t)}{D_k(t)}>\frac{D_k(t)}{D_{k+1}(t)}.
\]
因此序列 \((D_k(t))_{k\ge 0}\) 严格对数凸，而比值序列 \(\bigl(D_{k-1}(t)/D_k(t)\bigr)_{k\ge 1}\) 严格单调递减；
在边界反射坐标 \(q_k(t):=D_{k-1}(t)/D_k(t)\) 下，这与 Schur--Riccati 递推的单调收敛一致（推论 \ref{cor:pom-Lk-riccati-error-by-determinants}）。
\leanverified{detPoly\_eval\_pos}
\leanverified{detPoly\_eval\_ratio\_gap\_pos}
\leanverified{detPoly\_eval\_strict\_log\_convex}
\leanverified{detPoly\_eval\_strict\_mono}
\end{corollary}
\begin{proof}
由命题 \ref{prop:pom-Lk-det-cassini-pell}，当 \(t>0\) 时有 \(D_{k+1}D_{k-1}-D_k^2=t>0\)，即得严格对数凸性。
再将该恒等式除以 \(D_kD_{k+1}>0\) 得
\[
\frac{D_{k-1}}{D_k}-\frac{D_k}{D_{k+1}}=\frac{t}{D_kD_{k+1}}>0,
\]
从而比值严格单调。证毕。
\end{proof}
